% !Mode:: "TeX:UTF-8"

\cabstract{
针对管道输送系统中泄漏故障发现困难、定位精度不足的问题，本文围绕双通道时序传感信号，研究了一种基于深度学习的两阶段泄漏检测与定位方法。首先，依据原始 CSV 采样数据构建自动化数据处理流程，将长时序双通道信号按固定时间窗口切分为可训练样本，并分别生成用于分类任务和回归任务的数据清单。其次，针对泄漏识别与距离估计两个目标，设计了统一的一维卷积神经网络框架。其中，分类阶段以单通道分段信号为输入，实现正常状态与不同泄漏形态的识别；回归阶段以左右双通道信号对为输入，实现泄漏距离估计。最后，基于统一训练框架完成模型训练与评估，并采用准确率、平均绝对误差和均方根误差对方法性能进行分析。本文工作完成了从原始信号处理、任务定义到模型训练的完整流程，为后续开展正式实验、对比分析与工程化部署提供了基础。
}

\ckeywords{管道泄漏检测，时序信号分析，深度学习，一维卷积神经网络，距离回归}

\eabstract{
To address the difficulty of leak discovery and the insufficient localization accuracy in pipeline transport systems, this thesis studies a two-stage leak detection and localization method based on dual-channel time-series sensor signals. First, an automated data processing pipeline is constructed from raw CSV measurements. Long dual-channel signals are segmented with a fixed time window to generate trainable samples for both classification and regression tasks. Second, a unified one-dimensional convolutional neural network framework is designed for leak recognition and distance estimation. In the classification stage, segmented single-channel signals are used to identify normal states and different leak patterns. In the regression stage, paired left-right channel signals are used to estimate leak distance. Finally, a unified training framework is employed for model optimization and evaluation, and the method is analyzed with accuracy, mean absolute error, and root mean square error. The work in this thesis establishes a complete workflow from raw signal processing and task definition to model training, providing a foundation for subsequent formal experiments, comparative analysis, and engineering deployment.
}

\ekeywords{Pipeline Leak Detection, Time-series Signal Analysis, Deep Learning, One-dimensional Convolutional Neural Network, Distance Regression}

\makeabstract
